% Náhled k~sekci o~optimalizačních problémech: opakovaným dotazem na rozhodovací
% variantu najdeme optimum.
\begin{tikzpicture}[x=1cm, y=1cm]

    \tikzset{
        dotaz/.style={draw, thick, rounded corners=3pt, minimum width=30mm,
                      minimum height=8mm, align=center, font=\small}
    }

    \node[dotaz, fill=darkred!10]   (q1) at (0, 0)    {$\NZMNA(G,n)$?\quad \xmark};
    \node[dotaz, fill=darkred!10]   (q2) at (0,-1.2)  {$\NZMNA(G,n-1)$?\quad \xmark};
    \node                            (dots) at (0,-2.2) {$\vdots$};
    \node[dotaz, fill=darkgreen!15] (q3) at (0,-3.2)  {$\NZMNA(G,j)$?\quad \cmark};

    \draw[->, >=stealth, thick] (q1) -- (q2);
    \draw[->, >=stealth, thick] (q2) -- (dots);
    \draw[->, >=stealth, thick] (dots) -- (q3);

    \node[font=\small, darkgreen, anchor=west] at (2.6,-3.2)
        {největší nezávislá množina má $j$ vrcholů};

\end{tikzpicture}
